
\subsection{\acf{CbCC}}
Soll ein Programm von einem Computer ausgeführt werden,
ohne dass ihm vollständig vertraut werden kann,
erhält es zumindest Eingabeparameter
und ohne weitere Restriktionen hardware-spezifische und
ähnliche Informationen. Aus Sicht der Sicherheit besteht die
Möglichkeit, die Zugriffe des Programms auf absolut nötige
Ressourcen zu beschränken. Allerdings auch dann könnte das
Programm Informationen erhalten, die es zwar verarbeiten,
aber nicht an Dritte weiterleiten soll. Das Verhindern solcher
Informationsflüsse wurde bereits von \cite{lampson_note_1973}
als das „Confinement Problem“ beschrieben. Es bezieht sich auf
das Übertragen von Informationen eines Programms nicht an das
Programm oder den Nutzer, der es aufgerufen hat, sondern an den
„Owner“, also beispielsweise den Ersteller des Programms.
Die verschiedenen Wege der Informationsübertragung lassen sich
dabei in drei verschiedene Kategorien einordnen: in
Speicher, legitime Kanäle und sogenannte verdeckte, also
„Covert“-Channel, die implizit Informationen über nicht dafür
vorgesehene Zustände wie beispielsweise die Prozessorauslastung
übertragen. (\cite[613,614]{lampson_note_1973})

Wird auf diesem Wege eine Information nicht absichtlich von einem
„Insider“ an einen „Outsider“ gesendet, sondern kann dieser
Outsider aus der Beobachtung dieser Kanäle auf Informationen des
inneren Zustands eines Systems schließen, dann handelt es sich
um einen Seitenkanal, also einen Side-Channel.
(\cite[1]{almgren_c5_2015}, \cite[1]{wang_covert_2006})

Covert-Channel können grob in zwei Kategorien eingeteilt werden,
deren mögliche Verwendung stark vom jeweiligen Kontext abhängt.
\cite[3]{paccagnella_lord_2021} unterscheiden nach
eviction- und contention(Konkurrenz)-basierten Covert-Channels.
Evictionbasierte Channels bringen die Mikroarchitektur zunächst
in einen bekannten Zustand, um anschließend die Veränderung nach
der Ausführung des Senders zu beobachten und entsprechende
Informationen abzuleiten. Contentionbasierte Channels
encodieren dagegen die Information nicht in einem festen Zustand,
sondern in der Stärke der Konkurrenz um geteilte Ressourcen, die
sich beispielsweise als Zugriffszeit messen lässt.

Im Kontext von Transient Execution werden häufig
evictionbasierte Covert-Channel genutzt, da bei contentionbasierten 
Kanälen die Ausführung von Sender und Empfänger zeitgleich
stattfinden muss und damit nur ein binärer Wert je transient
Fenster übertragen werden kann. Dennoch existiert Forschung zu
diesen Kanälen, da viele Mitigationen gegen
klassische evictionbasierte Covert-Channel hier nicht wirken.
(\cite[8]{fustos_spectrerewind_2020})

\subsubsection{Cache als Informationsmedium}
Jedes komplexere Programm löst bei seiner Ausführung Prozesse in
der Speicherarchitektur aus. In höher abstrahierten Sprachen,
werden diese Prozesse nur implizit sichtbar, während bei Assembly
diese beispielsweise durch „mov“-Instruktionen ausgelöst
werden. Die x86-\ac{ISA} spezifiziert auch sehr granulare
Operationen auf der Mikroarchitektur des Speichers beispielsweise
durch die „prefetch“-Instruktion
(\cite{oracle_prefetch_2020}).

Der Zustand, ob eine bestimmte Adresse gecacht oder nicht gecacht
ist, kann als Information
interpretiert und damit als Covert-Channel betrachtet werden.
Da dieser Zustand nicht durch die \ac{ISA} sichtbar ist,
wird die Zugriffszeit auf eine Adresse gemessen und anhand eines
zeitlichen Schwellenwerts bestimmt, ob diese Adresse gecacht ist oder
nicht. Dieser Schwellenwert kann durch vorherige empirische
Messungen in demselben System festgelegt werden.
Neben der einfachen Encodierung von Werten durch den Zustand einer
Adresse im Cache existieren weitere Low-Level-Covert-Channels, die
ebenfalls den Cache oder Elemente der
Cachearchitektur nutzen. Ein solcher Ansatz wird in Anhang
\ref{appendix:advanced-cbcc} erläutert. In dieser Arbeit
liegt der Fokus jedoch auf Covert-Channels, bei denen
Informationen in den Zustand von Adressen kodiert werden.

Ursprünglich konzentrierten sich Angriffe mit \ac{CbCC} vor allem
auf Side-Channels zum Mitlesen kryptografischer Prozesse
zwischen Prozessen und virtuellen Maschinen
(\cite{bernstein_cache-timing_2005}, \cite{zhang_cross-vm_2012},
\cite{su_survey_2021}). Es wurden jedoch auch Covert-Channels mit
einem vom Angreifer kontrollierten Sender betrachtet, insbesondere
im Kontext von Cross-VM Kommunikation
(\cite{xu_exploration_2011}, \cite{maurice_hello_2017}).

Besonders relevant sind \ac{CbCC} seit
\cite{kocher_spectre_2019} jedoch für Angriffe mit
spekulativer Ausführung, wie in Abschnitt
\ref{section:spectre} erläutert wird.

\paragraph{Das Encodieren}
erfordert das Schreiben unabhängiger Zuständen in den Cache.
Dazu wird ein Speicherbereich
allokiert und einzelne Adressen werden als Informationswerte geladen
beziehungsweise evicted. Der Abstand zwischen diesen Adressen
wird als „Stride“ bezeichnet und muss mindestens der Größe eines
Cache-Blocks entsprechen, da sich sonst mehrere Informationsadressen
in einem Block befinden und  diese damit nicht mehr unabhängig sind.
Moderne \acp{CPU} implementieren einen Stride-Prefetcher, der Adressen mit
einem regelmäßigen Abstand innerhalb einer Memory-Page prefetchen kann
(\cite[739]{saileshwar_streamline_2021}). Aus diesem Grund sollte zur
Veringerung der Fehlerwahrscheinlichkeit als Stride die Größe
einer Memory-Page gewählt werden, um das Prefetchen von naheliegenden
Adressen zu verhindern. Diese Beobachtung zeigt sich in Abschnitt
\ref{section:experimental-results} an praktischen Messungen.

Für \ac{CbCC} wurden von \cite[13]{mcilroy_spectre_2019} zwei
Encodierungen betrachtet: „direct-mapped bits“ (bitweise) und
„indexed values“ (Array-Index). Bei der bitweisen Encodierung
symbolisiert jede verwendete Adresse ein Informationsbit durch
den Zustand, ob die Adresse gecacht oder nicht gecacht ist.
Bei der Array-Index-Encodierung wird die Information, als
Zahlenwert interpretiert und die Adresse mit diesem Index gecacht.
Das bedeutet in der Programmierung, dass $Array+(Stride*Secret)$
in den Cache geladen wird, wobei $Array$ die Startadresse des
Arrays ist.
Die Anzahl der während eines transient Fensters übertragenen
Informationen wird in dieser Arbeit in Bit betrachtet. Entsprechend
können während eines transient Fenster $n$ Bit übertragen werden,
wenn die Encodierung den Cache innerhalb des Fensters in $2^n$
verschiedene Zustände versetzen kann.

Der größte Vorteil von Array-Index-Encodierung besteht darin,
dass die Informationsübertragung grundsätzlich nicht von der
Länge des transient Fensters abhängt, da für die Encodierung
nur ein Memory Load oder Eviction nötig ist. Der Nachteil besteht
darin, dass der bitweise Informationsgehalt einer Übertragung
$\log_{2}\left(\frac{Array}{Stride}\right)$ beträgt und
$\frac{Array}{Stride}$ verwendete Adressen evictet und
zeitlich gemessen werden müssen. Bei der bitweisen Encodierung
müssen dagegen nur so viele Adressen
evictet und gemessen werden, wie Informationsbits übertragen
werden. Dafür ist die Anzahl der übertragenen Bits durch die
Größe des transient Fensters begrenzt. Eine gewisse
Kombination der beiden Konzepte ist auch möglich und wurde in
den \acp{POC} von Angriffen, wie beispielsweise
in \cite{hertogh_rain_nodate}, realisiert. Dort wurde ein Byte
in zweimal 16 Adressen, anstatt in 256 Adressen
encodiert. Dafür wurde das zu übertragende Byte in zwei 4-Bit
große Teile geteilt und der numerische Wert jeweils als Index
in ein Array mit 16 Adressen encodiert. Aus diesem Grund mussten
für diese Encodierung allerdings zwei Adressen während der
Transient Execution geladen werden, anstatt nur eine.

\paragraph{Die Übertragung}
durch evictionbasierte \ac{CbCC} erfolgt wie die meiste
Kommunikation in drei Schritten: Vorbereitung, Encodierung
und Auslesen. Die Vorbereitung besteht darin
das Medium in einen bekannten Zustand zu versetzen.
Die Encodierung verändert den Zustand des Mediums und bei dem
Auslesen wird der potenziell veränderte Zustand gemessen.
Während bei contentionbasierten Covert-Channels die Encodierung
und das Auslesen zeitgleich stattfinden, erfolgen diese Phasen
bei \ac{CbCC} sequenziell in der genannten Reihenfolge.
Im Kontext von Transient Execution können diese Phasen
sowohl in dem vom Angreifer kontrollierten, als auch im Ziel
Kontext stattfinden. Dadurch kann es passieren, dass die beiden
Prozesskontexte keine Speicherbereiche teilen und damit ein
direkter Zugriff auf geteilte Ressourcen nicht
möglich ist. (\cite[15]{xiong_survey_2022})

Um dennoch den Cache-Zustand einer Adresse
außerhalb des zugreifbaren Speichers zu ermitteln, werden
häufig Eviction-Sets verwendet. Dabei kann durch vorherige
Messungen bestimmt werden, welche der eigenen Adressen zu
einem Eviction-Set anderer Adressen im Ziel Kontext gehören
(\cite{vila_theory_2019}). Das Laden dieser Adressen im
Zielkontext führt zur Eviction einer Adresse im eigenen
Adressraum, wenn die physischen Einträge des
Eviction-Sets vorher gefüllt wurden
(\cite[607]{liu_last-level_2015}).

\paragraph{Störungen\label{section:cbcc-disturbance}} bei der
Übertragung durch \ac{CbCC} können
durch verschiedene mehr oder weniger beeinflussbare Faktoren
auftreten. Auch hier sind die Störungsquellen abhängig von der
konkreten Implementierung. Störungen und daraus folgende Fehler
wurden besonders bei contentionbasierten \ac{CbCC} untersucht,
wobei die Fehlerquellen nicht identisch zu den möglichen
Störungen bei evictionbasierten \ac{CbCC} sind. Bei \ac{CbCC}
ist ein zentrales Problem die Synchronisierung
und die damit verbundenen Timing-Probleme
(\cite[6,8]{maurice_hello_2017}),
was im Kontext von Transient Execution kein Problem darstellt.
Andere Störungsquellen sind bei Transient Execution dennoch
identisch zu contentionbasierten \ac{CbCC} zwischen zwei
verschiedenen Kontexten. Beispiele hierfür sind
das Laden anderer Adressen aus demselben Cache-Block durch
dasselbe Programm, zufälliges Laden durch den Prefetcher
der \ac{CPU} oder zufällige Evictions durch das Laden von
Adressen im selben Eviction-Set durch dasselbe oder ein anderes
Programm. Zusätzlich kann je nach \ac{CPU} Architektur das
Scheduling zu Fehlern führen. Wenn beispielsweise bei der Cache
Architektur aus Abbildung \ref{figure:cache-architecure-amd-a10}
der Prozess vor dem Auslesen auf einen anderen Kern verschoben
wird, ist es aufgrund des private \ac{LLC} nicht mehr möglich bei
dem Laden der Adresse den encodierten Zustand abzufragen.
(\cite[4]{wu_whispers_2012})
Covert-Channel, bei denen sich Sender und Empfänger nicht auf
demselben Kern befinden müssen, werden als Cross-Core-Channels
bezeichnet.

\subsubsection{Realisierung verschiedener Cache-Covert-Channel}
Obwohl das Kommunikationsmedium bei \ac{CbCC} grundsätzlich immer
der Zustand einer Adresse ist, also ob diese Adresse gecacht ist
oder nicht, kann sich die Methode für das Verändern des Zustands, das
Lesen des Zustands und auch die konkret gemessenen Zustände,
zwischen verschiedenen Techniken stark unterscheiden.

Viele einfache Cross-Core-\ac{CbCC} nutzen als Encodierung die
unterschiedlichen Zugriffszeiten, je nachdem, ob sich eine Adresse
im Cache oder nur im Speicher befindet. Die Methoden
zum Verändern beziehungsweise Lesen des Zustands umfassen dabei
meist die \verb|clflush|-Instruktion und das Laden
oder Schreiben von Adressen. (\cite[722]{yarom_flushreload_2014},
\cite[6]{almgren_c5_2015})

Diese Art von \ac{CbCC} wurde bei bisher demonstrierten
\acp{TEA} häufig verwendet (\cite[23]{xiong_survey_2022}), auch wenn
komplexere Techniken, die besonders die unteren Cacheschichten
verwenden, deutlich höhere Bandbreiten erreichen können,
teilweise auch bei Cross-Core Kommunikation und ohne geteilte
Speicherbereiche.
(\cite[13]{almgren_c5_2015}, \cite[10]{wang_bandwidthbreach_2023},
\cite[5]{rauscher_systematic_2025})

Wie bereits erwähnt, sind durch Eviction-Sets auch Covert-Channel
ohne geteilte Speicherbereiche möglich. Ist ein Eviction-Set
gefunden, kann durch gezieltes Laden dieser Adressen eine
Eviction ausgelöst und damit der Zustand des Caches in beiden
Speicherkontexten verändert werden. Der Angreifer kann anschließend
prüfen, ob eine seiner Adressen aus dem Eviction-Set aus dem Cache
entfernt wurde und dadurch ermitteln, welche Adresse geladen wurde.
Diese Technik wurde, wie viele andere \ac{CbCC}, ursprünglich als
Side-Channel Angriff auf Verschlüsselung angewendet
(\cite[10-13]{hutchison_cache_2006}) und später für Cross-Prozess
beziehungsweise Cross-VM Covert-Channel Kommunikation
(\cite[620]{liu_last-level_2015}) genutzt. Ein Nachteil dieses
Covert-Channels ist der Aufwand Eviction-Sets zu finden und diese
vor der Übertragung mit Blöcken zu füllen.

Welche \ac{CbCC} für einen Angriff verfügbar sind, hängt von der
Hardware, also der Anzahl und Eigenschaften der Cache-Ebenen, von der
\ac{ISA} sowie von dem Kontext des Ziel- und Angreiferprogramms ab. So
wird beispielsweise bei Flush+Reload die Adresse vom Zielkontext in
alle Cache-Ebenen geladen. Der Angreiferkontext kann beim Reload
die Adresse sowohl aus dem \ac{LLC}, als auch dem L1- oder L2-Cache
laden, je nachdem ob sich der Prozess auf demselben oder einem
anderen Prozessor Kern befindet (\cite[23]{xiong_survey_2022}) und
sofern von einem inclusive, private L1- und L2-Cache sowie einem
geteilten L3-Cache ausgegangen wird.

Schon durch kleine Änderungen können \ac{CbCC} je nach Kontext
optimiert werden. Ein Beispiel dafür ist Flush+Flush
als Abwandlung von Flush+Reload.
Anstatt eine Adresse zum Auslesen zu laden, wird die Adresse erneut
mit \verb|clflush| aus dem Cache entfernt. Wenn die Adresse gecacht
ist, weist diese Instruktion eine leicht höhere Laufzeit auf, die
gemessen werden kann. (\cite[4]{caballero_flushflush_2016})

Verschiedene Techniken werden in Anhang
\ref{appendix:cbcc-techniques} ausführlicher erläutert.
